\section{Signed compensation near a fixed double zero}
\label{sec:gt-critical-scale}

Suppose that the physical field \(F_T\) has an isolated double zero
at a real \(a\ne0\), with \(0<T<1/2\). Fix a circle of radius \(\rho\)
about \(a\) whose closed disk avoids \(0\) and lies in \(|\Im z|<1\).
Choose it small enough that \(F_T=(z-a)^2 A_T\) with \(A_T\)
nonvanishing there. By fixing a smaller concentric circle and
shrinking a compact time interval \(I\) about \(T\), both circles
have boundary floors and the physical pair stays in the smaller
disk. Set \(A_*=A_T(a)\ne0\) and \(Z_*=\max_\Gamma|z|\).
All limits here are for this fixed event.

\subsection{Signed event errors and the model double time}

Write \(p_j(t)=h_j^*(t)/(2j)!\) in this section and
\(\bar p_j=\sup_I p_j\). The phase bounds of
\eqref{eq:gt-physical-error} and \eqref{eq:gt-residual-phase-bound}
give \(E_{j,t}-1=O(\varepsilon_j)\),
\(\Psi_{j,t}-1=O(\varepsilon_j)\) on the fixed disk and time range,
where \(\varepsilon_j=\log^2(j+1)/j\). By
\eqref{eq:coefficient-envelope}, \(p_j\le B_{j-1}^j/(2j)!\),
so \(j^k\bar p_j M^{2j}\to0\) for any fixed \(M>0,k\ge0\).
The norm identity \eqref{eq:hflow}, with the normalization by \(Z_t\),
also gives
\begin{equation}
 (\log p_j)'=\beta_j+\beta_{j+1}-\beta_1,\qquad
 \left|\log\frac{p_j(t)}{p_j(T)}\right|
 \le (B_j+B_0)|t-T|.
 \label{eq:gt-relative-norm}
\end{equation}
In \eqref{eq:gt-relative-norm} the prime is the time derivative.

The field error \(O(p_j(t)Z_*^{2j})\) on \(\Gamma\) transfers
both counts for all large degrees. Uniform convergence on a fixed
slightly larger compact disk in the strip gives uniform bounds for
the model spatial derivatives by Cauchy's estimate. Common boundary
floors and an inner buffer then give \(K_j\to K\) uniformly by
Proposition~\ref{prop:gt-transfer}. The coordinate and field
estimates also give \(A_j-A=O(p_j(t)Z_*^{2j})\) on \(\Gamma\) by
division by \(Q_j,Q\), and then on the disk by maximum modulus.
The same boundary argument with the exact physical field at two
times shows \(A(t,\cdot)-A_T=O(|t-T|)\) on the disk.

Let \(e_j=D-D_j\), and write \(z_\pm=c\pm\sqrt D\) for the
unordered physical pair. At fixed \(t\), use the two moment
formulas with origin \(c(t)\). The first physical moment is zero.
The collar log of \(G_{j,t}/F_t\) equals
\(-\Delta_j/F_t+O(p_j(t)^2 Z_*^{4j})\) on \(\Gamma\), by the
positive physical floor and the logarithm series there. Integration
by parts gives \(c_j-c=O(p_j(t)Z_*^{2j})\) and
\[
 D_j-D=\left\langle (z-c)\frac{\Delta_j}{F_t}\right\rangle_\Gamma
               +O(p_j(t)^2Z_*^{4j}).
\]
In particular the square of the center difference belongs to the
quadratic remainder. Since
\((z-c)/Q=\tfrac12[(z-z_+)^{-1}+(z-z_-)^{-1}]\), including at a
double, this gives the first variation
\begin{equation}
 e_j(t)=-\frac12\sum_\pm\frac{\Delta_j(t,z_\pm)}{A(t,z_\pm)}
               +O(p_j(t)^2Z_*^{4j}).
 \label{eq:gt-gap-first-variation}
\end{equation}
The same half-sum identity in the exact defect
\eqref{eq:gt-model-gap} gives
\begin{equation}
 r_N(t)=-\frac12\sum_\pm\frac{R_N(t,z_{\pm,N})}{A_N(t,z_{\pm,N})}.
 \label{eq:gt-root-average-defect}
\end{equation}

\begin{proposition}[Fixed-event signed errors and uniqueness]
\label{prop:gt-fixed-event}
At the fixed event,
\begin{align}
 D_j(T)&\sim L_j(T,a)/A_* ,\notag\\
 r_N(T)&\sim -R_N(T,a)/A_* ,\qquad
 \frac{r_N(T)}{D_N(T)}
       \sim\frac{(2N)(2N-1)}{a^2},\qquad
 \frac{r_N(T)}{D_{N+1}(T)}
       \sim-\frac{\Pi_N}{\beta_{N+1}(T)a^4}.
 \label{eq:gt-fixed-ratios}
\end{align}
For all large \(N\), the model has one double time \(T_N\) on a
sufficiently small fixed interval about \(T\), and
\(T_N-T\sim-L_N(T,a)/(2A_*)\).
\end{proposition}
\begin{proof}
At \(T\) the physical roots in \eqref{eq:gt-gap-first-variation}
coincide at \(a\). Dividing the quadratic remainder by
\(p_j(T)|a|^{2j}\) leaves
\(O(p_j(T)(Z_*^2/|a|)^{2j})\), which tends to zero faster
than every inverse power of \(j\). Thus the asymptotic is relative
to the actual norm even when some links are small. For the residual
at \(T\), in \eqref{eq:gt-model-gap} replace the center and the
denominator by \(a,F_T\). The replacements cost
\(O(p_N(T)Z_*^{2N}\sup_\Gamma|R_N|)\), and the remaining simple
residue is \(-R_N(T,a)/A_*\). Division by
\(h_N^*(T)|a|^{2N-2}/(2N-2)!\) again leaves a factorial-small
relative remainder. Substituting the leading monomials proves the
ratios. The model residual and \(D_N(T)\) have eventual sign
\((-1)^N\operatorname{sgn}A_*\), and \(D_{N+1}(T)\) has its
opposite. There is no degree-independent favorable sign here.

The exact defect tends to zero uniformly on \(I\) by its contour
formula, fixed floor, and the factorial bound for
\(R_N=O(N^2\bar p_N Z_*^{2N-2})\). Together with \(D_N\to D\)
and \(K_N\to K\), the two gap laws give \(D_N'\to D'\) uniformly.
Since \(D(T)=0\), \(D'(T)=2\), choose the fixed interval so
\(D'\ge1\) and its endpoint signs are opposite. Then \(D_N\)
is strictly increasing with opposite endpoint signs for all large
\(N\). The mean value theorem gives
\(T_N-T=-D_N(T)/D_N'(\xi_N)\) with \(\xi_N\to T\), and
\(D_N'(\xi_N)\to2\). This proves the asserted displacement and
upgrades the existence statement of
Proposition~\ref{prop:double-collision-resolution} under the
additional uniform defect control.
\end{proof}

\subsection{The critical time scale}

For a bounded real \(s\), put \(v=T+s/N^2\) and
\[
 \begin{gathered}
 W_{N,v}(t)=\exp\!\left(4\int_t^v K_N(q)\,\dd q\right),\\
 \mathfrak R_N(s)=\int_T^v W_{N,v}r_N\,\dd t,\qquad
 \mathfrak H_N(s)=4\int_T^v W_{N,v}(K-K_N)D\,\dd t .
 \end{gathered}
\]
These are oriented integrals for \(s<0\). Let \(e_0=e_N(T)\) and
define the entire functions
\[
 \alpha=4/a^2,\qquad
 C_a(s)=\sum_{m\ge0}\frac{(8s/a^2)^m}{(2m)!},\qquad
 J_a(s)=\int_0^s C_a(q)\,\dd q,\qquad
 H_a(s)=C_a(s)-1-\alpha J_a(s).
\]

\begin{theorem}[Joint signed profiles at a fixed event]
\label{thm:gt-critical-profile}
For the fixed event and contour above, \(e_0\ne0\) for all
sufficiently large \(N\). Uniformly for \(s\) in a fixed compact
real interval,
\begin{align}
 \frac{e_N(v)}{e_0}&=C_a(s)+O(\varepsilon_N),&
 \frac{r_N(v)}{N^2e_0}&=-\alpha C_a(s)+O(\varepsilon_N),\notag\\
 \frac{\mathfrak R_N(s)}{e_0}&=-\alpha J_a(s)
                      +O(|s|\varepsilon_N),&
 \frac{\mathfrak H_N(s)}{e_0}&=H_a(s)+O(\varepsilon_N).
 \label{eq:gt-critical-profile}
\end{align}
The errors are absolute after the displayed normalization, including
at zero profile values. The signed combination is exactly
\begin{equation}
 e_N(v)+\mathfrak R_N(s)-\mathfrak H_N(s)=W_{N,v}(T)e_0 .
 \label{eq:gt-critical-combination}
\end{equation}
\end{theorem}
\begin{proof}
At \(t=T+s/N^2\), the law \eqref{eq:gt-gap-law} and bounded
coordinate derivatives give \(c(t)=a+O(|t-T|)\) and
\(D(t)=2(t-T)+O((t-T)^2)\). For real \(x,y\), with the square
root chosen nonnegative real or positive imaginary,
\(|\sqrt x-\sqrt y|\le\sqrt{|x-y|}\). Apply this to the
factorial-small coordinate error, and use the expansion of \(D\)
also at \(s=0\). For either the physical or model pair \(z_\pm^*\),
\[
 N(z_\pm^*-a)=\pm\sqrt{2s}+O(N^{-1}),\qquad
 k\log(z_\pm^*/a)=\pm 2\sqrt{2s}/a+O(N^{-1})
 \quad(k=2N,2N-2).
\]
The log here is the analytic branch near one. Thus the average of
\((z_\pm^*/a)^k\) is \(C_a(s)+O(N^{-1})\), with bounded terms
also when \(s<0\). The unit values at either pair are
\(A_*+O(N^{-1})\). The norm ratio is
\(1+O(\log^2(N+1)/N^2)\) by \eqref{eq:gt-relative-norm}.
For \(j=N,N+1\) the quadratic remainder of
\eqref{eq:gt-gap-first-variation}, divided by
\(p_j(T)|a|^{2j}\), remains factorial-small. Inserting the phase,
unit and power estimates into \eqref{eq:gt-gap-first-variation}
and \eqref{eq:gt-root-average-defect}, and dividing by the
event asymptotic of Proposition~\ref{prop:gt-fixed-event}, proves
the first two profiles. The factor of the second profile is
\(-\lim (2N)(2N-1)/(N^2a^2)=-\alpha\).

On the oriented interval \(W_{N,v}=1+O(|s|/N^2)\)
since \(K_N\) is uniformly bounded. Changing variables to
\(q=N^2(t-T)\) in the residual integral gives its profile with
the error \(O(|s|\varepsilon_N)\). Equation
\eqref{eq:gt-joint-error} gives the exact signed combination in
either orientation. It supplies the remaining integrated drift
profile without differentiating any time-dependent phase bound.
\end{proof}

For \(s>0\), put \(u=\sqrt{8s}/|a|\). Then
\(C_a=\cosh u\), \(\alpha J_a=u\sinh u-\cosh u+1\), and
\(H_a=2(\cosh u-1)-u\sinh u<0\). The last inequality follows
because \(u\sinh u-2(\cosh u-1)\) has value and first derivative
zero at zero and second derivative \(u\sinh u>0\).
For \(s<0\), put \(u=\sqrt{8|s|}/|a|\). The corresponding
values are \(\cos u\), \(1-\cos u-u\sin u\), and
\(2(\cos u-1)+u\sin u\). The limiting separate absolute total in
units of the signed total of \eqref{eq:gt-critical-combination} is
\(|C_a|+|\alpha J_a|+|H_a|\). On the positive side it is
\(1+2\alpha J_a\) and can grow arbitrarily large as the fixed
scaled endpoint increases. At \(s=-a^2\pi^2/8\) the three signed
terms are \(-1,-2,4\) in units of \(e_0\), with separate total
\(7\). Dividing by a profile at one of its zeros would require a
finer estimate.

\subsection{Direct calculation of the drift profile}
\label{sec:gt-drift-kernel}

There is also a spatial calculation of the drift difference.
Set \(k=2/a\), \(d^2=2s\), and \(Q_d(x)=x^2-d^2\). On two fixed
circles of radii \(\sqrt{2S}<R_0<R\), consider
\(Y_\zeta=Q_d-\zeta\e^{kx}\) for \(|s|\le S\).
For small complex \(\zeta\), the two counts transfer by Rouch\'e
on both circles. Let \(P_\zeta\) be the monic pair polynomial
and \(U_\zeta=Y_\zeta/P_\zeta\). Its coefficients, obtained by
contour moments, are holomorphic in \(\zeta\) through collisions.
The variation of the \(j\)-th moment at zero is
\[
 j\,\frac1{2\pi\ii}\oint_{|x|=R} x^{j-1}
          \frac{\e^{kx}}{Q_d(x)}\,\dd x .
\]
This follows by differentiating the logarithm and integrating by
parts. Inserting the first two moments into the pair polynomial
gives \(\partial_\zeta P_\zeta|_0=-I_d(x)\), where
\[
 I_d(x)=\cosh(kd)+x\,\frac{\sinh(kd)}{d}.
\]
This interpolates the exponential at \(\pm d\), with the Hermite
limit \(1+kx\) at \(d=0\). Set \(V_d=(\e^{kx}-I_d)/Q_d\), which
is holomorphic at the roots. Then
\(\partial_\zeta U_\zeta|_0=-V_d\). Since the unperturbed unit is
constant, the drift variation is
\(- (2\pi\ii)^{-1}\oint_{|x|=R} V_d'(x)/Q_d(x)\,\dd x\).
At \(d=0\) this gives \(-V_0''(0)=-k^4/12\). An
equivalent calculation from the full-field logarithmic derivative
in \eqref{eq:gt-regular-drift} is
\begin{align}
 \kappa_a(s)
   &=-\frac1{2\pi\ii}\oint_{|x|=R}
             \frac{2x\e^{kx}}{(x^2-2s)^3}\,\dd x\notag\\
   &=-\sum_{j\ge0}
       \frac{(j+1)(j+2)(2s)^j k^{2j+4}}{(2j+4)!}
    =\frac{C_a'(s)-k^2 C_a(s)}{8s}\quad(s\ne0).
 \label{eq:gt-drift-kernel}
\end{align}
Indeed the additional term from the denominator variation is
\(Q_d'I_d/Q_d^3\), rational with its poles inside and zero
contour integral by decay at infinity. Integration by parts
leaves the displayed integral. Expanding
\((1-2s/x^2)^{-3}\) gives the series. Comparing coefficients
with \(C_a(s)=\sum_{m\ge0}k^{2m}(2s)^m/(2m)!\) gives the last
identity. The series makes the extension entire and gives
\(\kappa_a(0)=-k^4/12=-4/(3a^4)\). It is negative for \(s\ge0\).
For \(s<0\), its expression is
\(k^4(\cos u-\sin u/u)/(4u^2)\), \(u=|k|\sqrt{-2s}\), and
it alternates in sign at positive integer multiples of \(\pi\).
For a polynomial perturbation of degree at most two, the
deflated factor is constant and the drift variation vanishes.
The interpolant above is therefore essential to the deflation.

To pass to the physical problem take \(x=N(z-a)\),
\(t=T+s/N^2\), and set
\[
 f_N(s,x)=N^2F_t(a+x/N)/A_* ,\qquad
 g_N(s,x)=N^2G_{N,t}(a+x/N)/A_* ,\qquad
 \zeta_N=N^2L_N(T,a)/A_*.
\]
The heat law gives \(\partial_tF_T(a)=-2A_*\).
Taylor expansion in space at \(T\), followed by integration of
the time derivative over the interval from \(T\) to \(t\), thus
gives \(f_N\to Q_d\) uniformly on compact \(x\)-disks and bounded
real \(s\)-intervals. By \eqref{eq:gt-physical-error},
\(g_N=f_N-\zeta_N B_N^{\rm sc}\), where
\[
 B_N^{\rm sc}(s,x)=
 \frac{h_N^*(t)}{h_N^*(T)}
       (1+x/(aN))^{2N}E_{N,t}(a+x/N)\longrightarrow\e^{kx}.
\]
The convergence follows from \eqref{eq:gt-relative-norm}, the
logarithm expansion of the monomial and the fixed-strip phase
bound. Also \(\zeta_N\to0\) by the factorial bound.

On a fixed sufficiently small complex \(\zeta\)-disk the
fields \(f_N-\zeta B_N^{\rm sc}\) are uniformly close to
\(Q_d\) on both contours, within their uniform boundary floors.
Their pair polynomials converge to those of \(Y_\zeta\) by the
moment formula. The denominators on the outer circle are bounded
below by \((R-R_0)^2\), and Cauchy's estimate on a slightly larger
disk in \(x\) gives uniform convergence of the spatial derivatives.
The drift integral consequently converges uniformly and is
holomorphic in \(\zeta\) on a common disk. Cauchy's formula there
gives convergence of the first parameter variation and a uniform
quadratic remainder. Hence the drift in the \(x\)-coordinate
changes by \(\zeta_N(\kappa_a(s)+o(1))\) between \(f_N\) and \(g_N\).
The counted roots in these auxiliary scaled disks belong to the
same fixed physical pair, by the counts on both boundaries.
A logarithmic derivative and a divided difference each contribute
\(1/N\) under the spatial rescaling, so \(K_z=N^2 K_x\) for
either field. Combining this with
\(-\zeta_N/(N^2e_0)\to1\) gives
\begin{equation}
 \frac{(K-K_N)(T+s/N^2)}{N^4e_0}\longrightarrow\kappa_a(s)
 \label{eq:gt-pointwise-transfer}
\end{equation}
uniformly on fixed compact real \(s\)-sets, through both mergers.
This derivation uses no time derivative of the approximate error.

\subsection{Shorter intervals and the adjacent head}

For any real nonzero \(\tau_N=o(N^{-2})\), put \(s_N=N^2\tau_N\).
By \eqref{eq:gt-critical-profile},
\[
 \mathfrak R_N(s_N)\sim\tau_N r_N(T),\qquad
 \mathfrak R_N(s_N)/e_0\sim-\alpha s_N .
\]
For \(q\) between \(0\) and \(s_N\), use
\(D(T+q/N^2)=2q/N^2+O(q^2/N^4)\) and
\(W_{N,v}=1+O(|s_N|/N^2)\).
The uniform additive kernel error in \eqref{eq:gt-pointwise-transfer}
then makes the normalized drift integrand in the scaled time
\(8q\kappa_a(0)+q\,o(1)\). Its oriented integral yields
\[
 \frac{\mathfrak H_N(s_N)}{e_0}\sim-\frac{16s_N^2}{3a^4},
 \qquad \frac{\mathfrak H_N(s_N)}{\mathfrak R_N(s_N)}
                  \sim\frac{4s_N}{3a^2}.
\]
This argument supplies a relative comparison even for intervals
shrinking faster than the uniform additive error.

The residual and the adjacent head share the first \(N\) link
indices. Set \(b_N^*=\beta_{N+1}(T)a^2/[(2N+1)(2N+2)]>0\).
The two actual event asymptotics and the physical leading ratio
\eqref{eq:leading-ratio} give
\[
 \frac{e_{N+1}(T)}{e_0}=-b_N^*[1+O(\varepsilon_N)],\qquad
 \frac{e_{N+1}(v)}{e_0}=-b_N^*[C_a(s)+O(\varepsilon_N)].
\]
For the second, apply Theorem~\ref{thm:gt-critical-profile} at
degree \(N+1\) and scaled argument \((1+1/N)^2s\).
No lower bound on the omitted link is needed. At the common
target \(T\) the unweighted combination in
\eqref{eq:gt-critical-combination}, divided by \(W_{N,v}(T)\),
is exactly \(e_0\). Its ratio in magnitude to the next-head error
is asymptotic to \(1/b_N^*\), at least of order
\(N^2/\log^2(N+1)\). The separately signed residual gives
\(b_N^*\mathfrak R_N(s)/e_{N+1}(T)
=\alpha J_a(s)+O(|s|\varepsilon_N)\). In particular,
\[
 \frac{\mathfrak R_N(s_N)}{e_{N+1}(T)}
     \sim\frac{\tau_N\Pi_N}{\beta_{N+1}(T)a^4}.
\]
At \(\tau_N=N^{-3}\) this is asymptotic to
\(16N/(\beta_{N+1}(T)a^4)\), at least of order
\(N/\log^2(N+1)\) in magnitude. In contrast, the absolute
fixed-circle allowance of Section~\ref{sec:gt-signed-cost},
using a fixed positive model floor \(m\), exceeds the
absolute signed residual by a factor asymptotic to
\(\rho(2\rho-b)|A_*|(Z_*/|a|)^{2N-2}/m\) on these shorter
forward slabs. This follows from \eqref{eq:gt-relative-norm},
the phase convergence, and
\(Q_{\tau_N}(M)/\tau_N\to1\) for a bounded drift allowance.
The statements compare a term, an explicit allowance, or an
exact matched error as specified. None bounds the cost of
acquiring a time slab from the cost of an endpoint snapshot.

The local analytic statements apply to each fixed source sharing
the stated norm and error hypotheses. The positive three- and
four-weighted-square controls have the same local event laws on
their own counted tubes. The score argument and the domination
used in \eqref{eq:gt-physical-error} hold for these finite square
sums on the half-line, with vanishing boundary term from the
factor \(u\) in the integration by parts. Their full theta
reflection and unit coefficients are absent, features unused
in the proof. The quadratic heat field
\((z-a)^2-2(t-T)\) gives \(D=2(t-T)\), \(K=0\) and exhausts
its incoming cost at \(T\), while lacking the infinite positive
Jacobi chain. The quantitative score and envelope hypotheses are
substantive. For example the Euler-factor \(c=1/2,p=23\)
control fails the stated \(B_0\) bound at time zero; this check
says nothing about arbitrary interior event intervals. The
reflected super-Gaussian fails the score confinement, and compact
and finite-atomic controls require separate envelope or chain
hypotheses. The canonical interior event considered here is
conditional. These fixed-tube comparisons do not supply a
whole-height collision exclusion or an exclusive advantage to
a transport algorithm.
